\documentclass[10pt]{article}
\usepackage[margin=20mm]{geometry}
\usepackage{amsmath,amssymb,mathtools,fontspec,unicode-math,xurl,hyperref}
\IfFontExistsTF{DejaVu Serif}{\setmainfont{DejaVu Serif}}{\setmainfont{Latin Modern Roman}}
\setmathfont{latinmodern-math.otf}
\setlength{\parindent}{0pt}\setlength{\parskip}{6pt}\sloppy
\title{Initial scientific review of M25-0026}
\author{}
\date{Review date: 25 October 2026}
\begin{document}
\maketitle
Abstract: No substantive issue was established in the inspected arguments of the manifest-pinned author manuscript. The review checks the approximation estimates, the prescribed itineraries and separation mechanism, the deduction of wandering Lakes of Wada, and the harmonic-measure argument. It does not certify every construction or the whole paper. Some geometric and external dependencies remain outside the completed verification.\par
\section*{Identity and scope}
The paper is David Martí-Pete, Lasse Rempe and James Waterman, Eremenko’s conjecture, wandering Lakes of Wada, and maverick points, arXiv:2108.10256v5, posted 30 October 2024, 43 PDF pages. The title page is dated 31 October 2024. The associated journal identifier is DOI 10.1090/jams/1049. All page references below are one-based pages of the acquired PDF; these agree with its printed page numbers. The original PDF SHA-256 is\par
> a835356f3b9ac20b05a86e666b1850c06828b16400473a8945966fc6614c7d80\par
All 25 manifest-listed files matched their recorded hashes. The journal final PDF was not acquired, and equivalence to that version has not been established. This review is confined to the author manuscript.\par
The extracted text of all 43 pages was read. Actual rendered PDF images were inspected for pages 1, 3, 6, 12–14, 20, 23–24, 27, 30–39. Reading all text does not mean that every assertion was independently proved. The principal completed checks are described below.\par
\section*{Definitions, quantifiers and dependencies}
For an entire function \(f\), the escaping set is \(I(f)=\{z:f^n(z)\to\infty\}\). The bungee set \(\operatorname{BU}(f)\) consists of points with unbounded orbits that do not tend to infinity. A compact set is full when its complement is connected. A wandering domain is a Fatou component whose forward component itinerary never repeats. Theorem 1.2 requires a nonempty, full, connected compact set. Theorem 1.7 allows a full compact set with countable prescribed sets \(Z_I,Z_{\operatorname{BU}}\), but no interior component may meet both. This restriction is necessary for the normal-family argument on a Fatou component; the paper does not prescribe incompatible behaviours within one component.\par
A maverick boundary point is distinguished by subsequence accumulation behaviour, not just by its unparameterised omega-limit set. Lemma 8.1 expresses this as \(\limsup_n\operatorname{dist}^\#(f^n(z),f^n(\zeta))>0\), where \(\zeta\) is an interior point. Boundedness of the domain itself is not boundedness of its orbits.\par
The core dependencies are the stated Arakelyan approximation theorem, uniform perturbation estimates in Section 2, and the strip system of Section 4. Section 7 does not depend on the boundary interpolation result of Section 9. Sections 5–6 use that result. The harmonic-measure argument uses Osborne–Sixsmith’s Lemma 4.1 and the previously established behaviour of multiply connected wandering domains.\par
\section*{Completed analytic checks}
\subsection*{Perturbations and limits}
In Lemma 2.3, pp. 12–13, fix \(\delta<\operatorname{dist}(A,\partial U)/2\) and \(\rho=\delta\inf_A|\phi'|/4\). Koebe’s theorem gives a uniform image disc of radius \(\rho\). For \(\varepsilon<\rho/2\), points outside the source \(\delta\)-disc cannot have the same perturbed image as its centre. Within that disc the two boundary curves differ by less than \(2\varepsilon<\rho\); the argument principle gives exactly one zero, counted with multiplicity. Cauchy’s estimate gives \(|f'-\phi'|\leq\varepsilon/\delta<\inf_A|\phi'|/2\). These arguments handle an unbounded closed \(A\) because the constants are uniform.\par
Lemma 2.5 explicitly requires a single modulus of continuity on the finite orbit union. With that interpretation, the recursive choice \[\delta_{n+1}(\epsilon)=
 \min\{\delta_n(\epsilon),\delta_n(\delta'(\epsilon)),\epsilon/2\}\] keeps the perturbed orbit in the domain and bounds the next error by \(\epsilon\). For Corollary 2.7 one may first thicken \(K\) inside \(G\); the estimates on this neighbourhood supply the uniform closeness needed for the injectivity argument. This is the relevant reading of its use in the unbounded construction.\par
The approximation sets in the inductive constructions exhaust the plane and preserve previous sets. If \(\epsilon_{j+1}\leq\epsilon_j/2\), then the tail error relative to \(g_j\) is at most \(2\epsilon_j\). The stability thresholds are chosen for this total error before approximation. Consequently the locally uniform entire limit retains the finitely many requirements imposed at each stage. Mere convergence on compact sets, without this uniform tail control on the unbounded sets, would not suffice.\par
\subsection*{Strip itinerary and prescribed escaping behaviour}
The arithmetic in (4.1)–(4.2), p. 20, gives \[N_j=\frac{j(j+3)}2,\qquad N_{j+1}-N_j=j+2.\] For \(\Phi(z)=5z\), the image of \(S_j\) contains both \(S_{j+1}\) and \(T_{j+1}\) with positive vertical margins. If \(|f-\Phi|<1/2\), then \[|\operatorname{Re}f(z)|\geq5|\operatorname{Re}z|-1/2
   >4|\operatorname{Re}z|\quad (|\operatorname{Re}z|\geq1).\] Thus the expansion applies on either side of the imaginary axis.\par
The conditions of Proposition 5.2 imply unbounded orbits for all points of \(K\) by visits to \(T_j\). A prescribed bungee point has infinitely many visits to the bounded part of \(S_0\) because it occurs infinitely often in the scheduling sequence. A prescribed escaping point eventually belongs to every finite set \(Z_j\), and condition (e) gives a lower bound \(j\) throughout each interval \([N_j+1,N_{j+1}]\). These intervals cover all sufficiently large times, not merely a subsequence.\par
There is also a direct check of pairwise disjointness of the compacta. If \(f^n(K)\) and \(f^m(K)\) intersect with \(n<m\), choose \(j\) sufficiently large that \(N_j>m\) and \(j+1>m-n\). Iterating an intersection point for \(N_j-m\) steps would give an intersection of \(T_j\) with an image in one of the preceding \(S\) strips. This contradicts their disjointness. Interior normality follows because all images omit a fixed open disc; the approximating boundary sets entering that disc show \(\partial K\subseteq J(f)\). Hence the interior components are whole Fatou components, with wandering itineraries.\par
Lemma 5.4, p. 24, uses simple connectedness and the lower bound for the hyperbolic metric to obtain \[\operatorname{dist}_{U_m^\zeta}(\zeta,\omega)
 \geq\frac12\log\left(1+\frac{|\zeta-\omega|}{2\delta_m}\right)
 \longrightarrow\infty.\] The crossing of \(\partial K\) is justified both when an endpoint is on that boundary and when the endpoints are in different interior components. Countability of the prescribed points allows the subsequent diagonal subsequence selection.\par
\subsection*{Connected components, rather than only paths}
For Theorem 7.1, pp. 31–35, the half-strip boundaries enter a bounded forward-invariant disc. They are disjoint from \(J(f)\cup I(f)\cup\operatorname{BU}(f)\). For every point outside \(K=\bigcap_jK_j\), some such boundary separates it from \(K\). This excludes a connected subset of that union from meeting both \(K\) and its complement; the argument does not merely exclude paths.\par
For each fixed \(t>0\), the inclusion \(1/(j+2)\leq t\leq j+2\) holds for all sufficiently large \(j\). Thus condition (e) gives bounded returns for every non-endpoint of the ray, while the visits to \(T_j\) give unbounded suborbits. Condition (f), on the full intervening time intervals, makes the endpoint escape. These facts imply that the endpoint is a singleton component of \(I(f)\). The bungee orbit also rules out a polynomial of degree at least two; the prescribed strip expansion and bounded invariant disc rule out the affine cases.\par
Lemma 7.2’s derivative estimate is consistent with the change of variables \(\zeta=e^{-z}\) and \(\psi(\zeta)=\lambda\zeta+O(\zeta^2)\): \[\phi'(z)=\frac{\zeta\psi'(\zeta)}{\psi(\zeta)}
          =1+O(e^{-\operatorname{Re}z}).\] This supplies control at the unbounded end; compactness and univalence supply bounds on the remaining, strictly interior portion. In the general-continuum construction, equation (7.3) has an absolute value, verified in both the actual PDF image and source line 2023. The plain-text extraction dropped this symbol. A purported sign objection based on that extraction would be invalid.\par
\subsection*{Lakes of Wada and harmonic measure}
Given a Lakes of Wada continuum \(X\), take the complement of its unbounded complementary domain to obtain the full compactum used in Theorem 1.7. Its bounded complementary domains are interior components with boundary \(X\). Prescribing escaping behaviour in one and oscillating behaviour in another is permitted. Harmonic measures seen from these different domains need not agree. Thus having full harmonic measure for escaping points from one domain and for nonescaping points from another is not a contradiction.\par
For Theorem 1.11, pp. 37–38, select precisely those times when the interior orbit lies within spherical distance \(\delta\) of \(\xi\). The maps between those successive wandering domains are iterates of \(f\), and their compositions are the selected iterates. This ensures the eventual interior proximity assumption of the original Osborne–Sixsmith Lemma 4.1, whose statement was checked in its primary published PDF, p. 567. The change from all sufficiently large times in that lemma to selected times in the present paper is justified by this reindexing.\par
For a maverick point, compactness of the sphere gives a subsequence of interior images tending to \(\xi'\) with boundary images a definite distance away. Choose rational \(0<\delta<1\) sufficiently small, and \(\xi\in\mathbb Q(i)\) within \(\delta/2\) of \(\xi'\). Triangle inequalities give interior distance less than \(\delta\) and boundary distance greater than \(2\delta\) infinitely often. The countable union over such \(\xi,\delta\) therefore covers every maverick point and has harmonic measure zero. Rational finite points are dense in the sphere, so the case \(\xi'=\infty\) is included.\par
Finally, for the finite-set formula (9.1), p. 39, \(\operatorname{Im}\psi'(z)>0\) throughout the upper half-plane. For two distinct points, integrate \(\psi'\) on their straight segment: the quotient \((\psi(z_2)-\psi(z_1))/(z_2-z_1)\) still has positive imaginary part. This proves global injectivity, rather than only nonvanishing derivative. The integral version has the same positivity, conditional on the stated Evans measure theorem.\par
\section*{Limits, alternative readings and impact}
No substantive Problem is asserted. Several apparent transcription issues do not establish a failure: (7.3) retains absolute values, and in the disc example on p. 30 the needed convergence is \(\pi_m(Z_{\rm mav})=(r/r_m)\Xi\to\Xi\), regardless of the displayed equality there. Typographical errors alone are outside the assigned scope.\par
The following work remains: a full independent proof of the uniform kernel-convergence assertion for arbitrary full continua in Theorem 1.2; the positive-area Jordan-arc construction and its conformal collapse in Theorem 1.14; reconstruction of the omitted construction in the second part of Proposition 7.6; verification of every auxiliary remark, including the fast-escape variant; and full proofs of the external approximation, potential-theoretic and topological results. No computational experiment is needed to replace these analytic tasks.\par
The consulted arXiv record identifies v5 as the author accepted manuscript and notes restructuring of Section 9. No later revision or explicit correction was visible on that record. The attempted official journal-page lookup returned HTTP 403. This is not an exhaustive errata search and does not prove journal-version equivalence.\par
\section*{Conclusion}
Item | Recorded outcome\par
Scientific judgment | No Issue Found within the completed checks\par
Core impact | No supported invalid derivation or counterexample established\par
Repair verification | No substantive repair proposed\par
Human acceptance | Pending; not filled on the user’s behalf\par
Author adoption | Not requested\par
Remaining science | Geometric and external dependencies listed above\par
Report delivery | TeX, compiled PDF, build and page-inspection evidence saved\par
9 D. Martí-Pete, L. Rempe and J. Waterman, Eremenko’s conjecture, wandering Lakes of Wada, and maverick points, arXiv:2108.10256v5 (2024). https://arxiv.org/abs/2108.10256v5. J. W. Osborne and D. J. Sixsmith, On the set where the iterates of an entire function are neither escaping nor bounded, Ann. Acad. Sci. Fenn. Math. 41 (2016), 561–578, DOI 10.5186/aasfm.2016.4134. Primary PDF statement of Lemma 4.1 consulted: https://www.acadsci.fi/mathematica/Vol41/vol41pp561-578.pdf.\par
\section*{Supporting checks}
Their evidence directory contains the input validation, extracted page text, rendered original-PDF pages, source fragment, explicit visual-inspection ledger and compilation-status record. The source inspected directly was `Eremenkos\_Conjecture\_v32.tex`, lines 1992–2055, SHA-256\par
> 58ab560be8a3f311c8adc55e7b3778d52c70de165a032f72f0255d3f9dbba417.\par
All analytic checks and actual page-image inspections were performed by the current assistant. No author code was run or imported. No numerical computation is claimed. Poppler rendered the original PDF; this is not compilation of the review report.\par
The authorized followup uses the existing Tectonic 0.17.0 runtime. Exact-source build logs, input and PDF hashes, and report-page inspection records are retained in the accompanying provenance and followup evidence. No compiler was installed by this conversation. The source PDF’s first-page rendering emitted a Symbol display-font warning; the displayed page was inspected and its title/version and reviewed mathematical content were legible.\par
No independent backend-route evidence or model-alias equivalence is asserted. No dispatcher tab was inspected or taken over.\par
\end{document}
